% ============================================================
\section{Coherence Gap: Spectral Lower Bounds from PSD Defect}
\label{sec:coherence-gap}
% ============================================================

\begin{quote}\small
\textbf{What this automates.}
When the defect energy operator $\CCDefOp{d}{h} = \CCLeak{d}{h}^* \CCLeak{d}{h}$ is strictly positive on the observable subspace, this positivity translates into a spectral gap.
This section provides the rigorous mechanism for converting ``defect is strictly positive'' into a quantitative lower bound on spectra.
\end{quote}

\subsection{Defect energy operator}

Recall from the Horizon Fundamental Theorem (Section~\ref{sec:hft}) that the leakage operator is $\CCLeak{d}{h} = \CCLeakExpr{d}{h}$ and the defect energy is $\CCDefEnergy{h}{x} = \|\CCLeak{d}{h} x\|^2$.

\begin{definition}[Defect energy operator]
\label{def:defect-energy-gap-operator}\lean{defectEnergyOp}
The \emph{defect energy operator} is
\[
\CCDefOp{d}{h} := \CCLeak{d}{h}^* \CCLeak{d}{h} = \CCPi{h} d^* \CCPibar{h} d \CCPi{h}.
\]
This is a positive semidefinite operator on $\cH$ with $\langle x, \CCDefOp{d}{h} x \rangle = \|\CCLeak{d}{h} x\|^2 = \CCDefEnergy{h}{x}$.
\end{definition}

\paragraph{PSD structure.}
Since $\CCDefOp{d}{h} = \CCLeak{d}{h}^* \CCLeak{d}{h}$, we have $\CCDefOp{d}{h} \succeq 0$ automatically.
The question is whether $\CCDefOp{d}{h}$ has a \emph{gap}, that is, a strictly
positive lower bound on the observable subspace.

\subsection{Spectral gap from PSD lower bound}

\begin{theorem}[Coherence gap theorem]
\label{thm:coherence-gap}\lean{coherence_gap_theorem}
Let $\CCDefOp{d}{h} = \CCLeak{d}{h}^* \CCLeak{d}{h} \succeq 0$ be the defect energy operator.
Suppose there exists $\mu > 0$ such that
\[
\CCDefOp{d}{h} \succeq \mu \, \CCPi{h}
\]
as operators on $\cH$, meaning $\langle x, \CCDefOp{d}{h} x \rangle \ge \mu \|\CCPi{h} x\|^2$ for all $x \in \cH$.

Then for all $x \in \cH_{\mathrm{obs}} = \CCPi{h} \cH$:
\[
\langle x, \CCDefOp{d}{h} x \rangle \ge \mu \|x\|^2,
\]
and consequently the spectrum of $\CCDefOp{d}{h}$ restricted to $\cH_{\mathrm{obs}}$ satisfies
\[
\mathrm{spec}(\CCDefOp{d}{h}|_{\cH_{\mathrm{obs}}}) \subseteq [\mu, \infty).
\]
\end{theorem}

\begin{proof}
For $x \in \cH_{\mathrm{obs}}$, we have $\CCPi{h} x = x$, so $\|\CCPi{h} x\|^2 = \|x\|^2$.
The hypothesis $\CCDefOp{d}{h} \succeq \mu \CCPi{h}$ gives
\[
\langle x, \CCDefOp{d}{h} x \rangle \ge \mu \|\CCPi{h} x\|^2 = \mu \|x\|^2.
\]
By the Rayleigh quotient characterization of spectra, if $\langle x, \CCDefOp{d}{h} x \rangle \ge \mu \|x\|^2$ for all nonzero $x$ in a subspace, then all eigenvalues of $\CCDefOp{d}{h}$ on that subspace are at least $\mu$.
\end{proof}

\begin{quote}\small
\textbf{Intuition.}
The coherence gap theorem converts a uniform positivity condition on defect energy into a spectral statement.
If every observable state leaks at least $\mu$ units of energy into the shadow, then no observable state can have defect energy below $\mu$.
This is the calculus mechanism for ``the defect cannot be arbitrarily small on observable states.''
\end{quote}

\subsection{Converse and gap detection}

\begin{lemma}[Gap detection via Rayleigh quotient]
\label{lem:gap-detection}\lean{gap_detection_ratio}
The largest $\mu \ge 0$ such that $\CCDefOp{d}{h} \succeq \mu \CCPi{h}$ is
\[
\mu_{\min} := \inf_{x \in \cH_{\mathrm{obs}},\, x \ne 0} \frac{\langle x, \CCDefOp{d}{h} x \rangle}{\|x\|^2}
= \inf_{x \in \cH_{\mathrm{obs}},\, x \ne 0} \frac{\|\CCLeak{d}{h} x\|^2}{\|x\|^2}.
\]
If $\mu_{\min} > 0$, then $\CCDefOp{d}{h}$ has a coherence gap of size $\mu_{\min}$ on $\cH_{\mathrm{obs}}$.
\end{lemma}

\begin{proof}
This is the standard Rayleigh quotient characterization of the smallest eigenvalue (or infimum of the spectrum) of a self-adjoint operator restricted to a subspace.
\end{proof}

\paragraph{Zero gap.}
If $\mu_{\min} = 0$, then there exist observable states with arbitrarily small
defect energy, so the defect can be screened by choosing appropriate states.
The coherence gap theorem does not apply in this case.

\begin{corollary}[Kernel and gap]
\label{cor:kernel-gap}\lean{kernel_gap_injective_statement}
If $\Ker(\CCLeak{d}{h}) \cap \cH_{\mathrm{obs}} = \{0\}$, then $\CCDefOp{d}{h}|_{\cH_{\mathrm{obs}}}$ is injective.
If additionally $\cH_{\mathrm{obs}}$ is finite-dimensional, then $\CCDefOp{d}{h}|_{\cH_{\mathrm{obs}}}$ has a strictly positive coherence gap.
\end{corollary}

\begin{proof}
Injectivity follows from $\langle x, \CCDefOp{d}{h} x \rangle = \|\CCLeak{d}{h} x\|^2 = 0 \Rightarrow \CCLeak{d}{h} x = 0 \Rightarrow x = 0$ for $x \in \cH_{\mathrm{obs}}$.
In finite dimensions, a self-adjoint positive definite operator has strictly positive smallest eigenvalue.
\end{proof}

\subsection{Finite shadow-transport invariants}

On a finite observed sector, the defect operator carries more than a single gap.
It determines a complete finite invariant family: spectrum, trace moments, and a
determinant polynomial. These are exact structural invariants of the horizon
cut, not additional hypotheses.

\begin{definition}[Shadow-transport invariant family]
\label{def:shadow-transport-invariants}\lean{shadow_transport_invariants}
Assume $\cH_{\mathrm{obs}}=\CCPi{h}\cH$ is finite-dimensional and define
\[
T_h := \CCDefOp{d}{h}\big|_{\cH_{\mathrm{obs}}}
=
\CCLeak{d}{h}^* \CCLeak{d}{h}\big|_{\cH_{\mathrm{obs}}}.
\]
The associated \emph{shadow-transport invariant family} at horizon $h$ consists
of:
\begin{enumerate}[label=(\roman*), itemsep=0.2em]
\item the multiset
      \[
      \Lambda_h(d):=\mathrm{spec}(T_h)
      \]
      of nonnegative eigenvalues of $T_h$;
\item the spectral moments
      \[
      M_k(h):=\Tr(T_h^k)
      \qquad (k\ge 1);
      \]
\item the determinant polynomial
      \[
      \Delta_h(t):=\det(I+tT_h).
      \]
\end{enumerate}
\end{definition}

\begin{theorem}[Finite shadow-transport package]
\label{thm:shadow-transport-package}\lean{shadow_transport_package}
Assume $\cH_{\mathrm{obs}}$ is finite-dimensional and let
$r:=\dim(\cH_{\mathrm{obs}})$. Then:
\begin{enumerate}[label=(\roman*), itemsep=0.2em]
\item every element of $\Lambda_h(d)$ is nonnegative;
\item the following are equivalent:
      \[
      \Lambda_h(d)=\{0,\dots,0\},
      \qquad
      T_h=0,
      \qquad
      \CCLeak{d}{h}\big|_{\cH_{\mathrm{obs}}}=0;
      \]
\item if $U$ is unitary and preserves both $d$ and $\CCPi{h}$, then
      $\Lambda_h(d)$, all moments $M_k(h)$, and $\Delta_h(t)$ are unchanged;
\item if $\Lambda_h(d)=\{\lambda_1,\dots,\lambda_r\}$ with multiplicity, then
      \[
      M_k(h)=\sum_{j=1}^r \lambda_j^k,
      \qquad
      \Delta_h(t)=\prod_{j=1}^r (1+t\lambda_j).
      \]
      Hence either the determinant polynomial $\Delta_h$ or the finite list of
      moments $M_1(h),\dots,M_r(h)$ determines the full multiset
      $\Lambda_h(d)$.
\end{enumerate}
\end{theorem}

(Proof in Appendix~\ref{sec:appendix-elimination}.)

\begin{corollary}[Symmetry factorization of the invariant family]
\label{cor:shadow-transport-factorization}\lean{shadow_transport_factorization}
Assume $\cH_{\mathrm{obs}}$ is finite-dimensional and carries a unitary action
of a finite group $G$ preserving both $d$ and $\CCPi{h}$. Let
\[
\cH_{\mathrm{obs}}=\bigoplus_\rho \cH_\rho
\]
be the isotypic decomposition from
Theorem~\ref{thm:symmetry-reduction-admissibility}, and write
$T_{h,\rho}:=T_h|_{\cH_\rho}$. Then
\[
T_h=\bigoplus_\rho T_{h,\rho},
\qquad
M_k(h)=\sum_\rho \Tr(T_{h,\rho}^k),
\qquad
\Delta_h(t)=\prod_\rho \det(I+tT_{h,\rho}).
\]
Thus the shadow-transport invariant family decomposes into independent symmetry
blocks.
\end{corollary}

\begin{proof}
By symmetry reduction, $T_h$ preserves each isotypic component and is
block-diagonal across the decomposition $\cH_{\mathrm{obs}}=\bigoplus_\rho
\cH_\rho$. Trace is additive on block sums and determinant is multiplicative on
block-diagonal operators, giving the stated formulas.
\end{proof}

\paragraph{Two-mode reductions.}
\label{rem:two-mode-reduction}
In a two-dimensional reduced block, the determinant and trace already fix the
entire spectrum. If $T$ is such a block, then its eigenvalues are the two roots
of
\[
\lambda^2-\Tr(T)\lambda+\det(T)=0.
\]
Thus whenever admissibility or symmetry reduction cuts a candidate family down
to a rank-two block with fixed trace and determinant data, the remaining choice
is an explicit algebraic branch-selection problem. The old trace/determinant
``trick'' is therefore not discarded; it appears here as the rank-two edge of a
general shadow-transport invariant family.
